\begin{table}[!ht]
\centering
\caption{\label{tab:desc_stat}Descriptive Statistics}
\centering
\resizebox{\ifdim\width>\linewidth\linewidth\else\width\fi}{!}{
\begin{threeparttable}
\begin{tabular}[t]{lcccc}
\toprule
\multicolumn{1}{c}{} & \multicolumn{1}{c}{\thead{Full sample\\\footnotesize{\textit{2009-2014 Bac cohorts}}}} & \multicolumn{1}{c}{\thead{Need-based grant\eligibles\\\footnotesize{\textit{in HS grad. year}}}} & \multicolumn{1}{c}{\thead{Merit aid eligibles\\\footnotesize{\textit{in HS grad. year}}\\\footnotesize{\textit{Bac grade $\geq 16$}}}} & \multicolumn{1}{c}{\thead{RD sample\\\footnotesize{\textit{Bac grade $\in (15, 17)$}}}} \\
 & (1) & (2) & (3) & (4)\\
\midrule
\addlinespace[0.3em]
\multicolumn{5}{l}{\textit{Panel A. Socio-demographic characteristics}}\\
\hspace{1em}Female (\%) & 52.3 & 56.6 & 59.3 & 58.5\\
\addlinespace[\smallskipamount]
\hspace{1em}Age at Bac (mean) & 18.8 & 18.5 & 18.0 & 18.1\\
\addlinespace[\smallskipamount]
\hspace{1em}Parents' taxable income (median; euros) & - & 21,492 & 28,910 & 27,133\\
\addlinespace[\smallskipamount]
\hspace{1em}Very high SES (\%) & 28.6 & 14.4 & 31.9 & 26.2\\
\addlinespace[\smallskipamount]
\hspace{1em}High or middle SES (\%) & 38.5 & 43.2 & 44.6 & 45.6\\
\addlinespace[\smallskipamount]
\hspace{1em}Low SES (\%) & 26.6 & 37.3 & 20.8 & 25.0\\
\addlinespace[\smallskipamount]
\hspace{1em}Missing SES (\%) & 6.3 & 5.1 & 2.6 & 3.2\\
\addlinespace[\smallskipamount]
\addlinespace[0.3em]
\multicolumn{5}{l}{\textit{Panel B. Academic characteristics}}\\
\hspace{1em}>= 16/20 at Bac (\%) & 5.5 & 5.0 & 100.0 & 47.7\\
\addlinespace[\smallskipamount]
\hspace{1em}Bac grade (mean) & 12.1 & 12.0 & 16.8 & 15.8\\
\addlinespace[\smallskipamount]
\hspace{1em}General high school track (\%) & 53.4 & 59.5 & 88.9 & 80.4\\
\addlinespace[\smallskipamount]
\hspace{1em}Technologic high school track (\%) & 22.3 & 26.5 & 7.5 & 11.6\\
\addlinespace[\smallskipamount]
\hspace{1em}Professional high school track (\%) & 24.3 & 14.0 & 3.6 & 8.0\\
\addlinespace[\smallskipamount]
\hspace{1em}Private high school (\%) & 21.9 & 15.0 & 21.0 & 20.0\\
\addlinespace[\smallskipamount]
\addlinespace[0.3em]
\multicolumn{5}{l}{\textit{Panel C. Financial aid}}\\
\hspace{1em}Eligible for merit aid in HS grad. year & 55,347 & 55,347 & 55,347 & 35,879\\
\addlinespace[\smallskipamount]
\hspace{1em}Eligible for merit aid in HS grad. year + 1 & 50,318 & 49,122 & 49,122 & 31,767\\
\addlinespace[\smallskipamount]
\hspace{1em}Eligible for merit aid in HS grad. year + 2 & 43,658 & 42,754 & 42,754 & 27,431\\
\addlinespace[\smallskipamount]
\hspace{1em}Echelon 0-0 bis (\%) & 7.6 & 23.1 & 39.1 & 35.1\\
\addlinespace[\smallskipamount]
\hspace{1em}Echelon 1 (\%) & 5.4 & 16.5 & 18.9 & 18.8\\
\addlinespace[\smallskipamount]
\hspace{1em}Echelon 2-4 (\%) & 8.5 & 25.9 & 22.0 & 23.5\\
\addlinespace[\smallskipamount]
\hspace{1em}Echelon 5-7 (\%) & 11.3 & 34.5 & 20.0 & 22.6\\
\addlinespace[\smallskipamount]
\addlinespace[0.3em]
\multicolumn{5}{l}{\textit{Panel D. Higher education outcomes}}\\
\hspace{1em}Enrolled in HS grad. year (\%) & 63.4 & 90.4 & 95.7 & 94.5\\
\addlinespace[\smallskipamount]
\addlinespace[0.3em]
\multicolumn{5}{l}{\textit{\footnotesize Among enrolled:}}\\
\hspace{1em}\hspace{1em}Public university (\%) & 48.3 & 52.4 & 36.6 & 39.7\\
\addlinespace[\smallskipamount]
\hspace{1em}\hspace{1em}Vocational degree (BTS) (\%) & 21.2 & 23.3 & 5.1 & 10.1\\
\addlinespace[\smallskipamount]
\hspace{1em}\hspace{1em}Technical degree (DUT) (\%) & 12.4 & 12.8 & 5.7 & 9.6\\
\addlinespace[\smallskipamount]
\hspace{1em}\hspace{1em}Academic preparatory classes (CPGE) (\%) & 9.7 & 6.8 & 37.5 & 28.8\\
\addlinespace[\smallskipamount]
\hspace{1em}\hspace{1em}Other institutions (\%) & 8.4 & 4.7 & 15.0 & 11.7\\
\addlinespace[\smallskipamount]
\hspace{1em}Enrolled in 2nd year in HS grad. year + 1 (\%) & 36.4 & 49.4 & 78.4 & 73.7\\
\addlinespace[\smallskipamount]
\hspace{1em}Enrolled in 3rd year in HS grad. year + 2 (\%) & 22.6 & 28.1 & 66.2 & 58.2\\
\addlinespace[\smallskipamount]
\hspace{1em}Obtained a degree (2009-2023) (\%) & 41.0 & 50.8 & 85.0 & 80.1\\
\midrule
\addlinespace[\smallskipamount]
Observations & 3,363,714 & 1,101,658 & 55,347 & 75,188\\
\bottomrule
\end{tabular}
\begin{tablenotes}[para]
\item \textit{Notes:} This table shows descriptive statistics for four samples: (1) the full sample of 2009-2014 high school graduates, (2) the analytical sample, i.e., students from these cohorts who are eligible to a need-based grant in their high school graduation year, (3) students eligible for the merit aid in that same year, i.e., students from the analytical sample who obtained at least 16/20 at the Bac, and (4) the RD sample, i.e., students from the analytical sample who obtained between 15 and 17 at the Bac. HS grad. stands for high school graduation. Since about 5\% of students in the full sample have multiple enrollments in their high school graduation year, I follow \citet{bonneau_quelle_2021}'s ranking across enrollments to assign them a main enrollment. This ranking is based on knowledge of the French higher education system. Degree completion is not observed for vocational degrees (BTS) and some other programs; students initially enrolled in these programs are therefore excluded from this statistic. The divergence in merit-aid-eligible counts between columns (1) and (2) in years +1 and +2 reflects students who became eligible for a need-based grant only after enrolling in higher education.
\end{tablenotes}
\end{threeparttable}}
\end{table}
